% !TEX root = userguide.tex

\def\headdate{15th June 2009}

\section{Iterative solvers for Stokes and Navier-Stokes}
\label{sec:iterativestokes}

In this section some examples using the block preconditioned iterative solvers
for the Stokes and Navier-Stokes equations as described in the book
by Elman \emph{et al.} \cite{Elman05} are given.

\subsection{Theory}

The system of equations that needs to be solved after discretization in time
and space of the Navier-Stokes equations with constant 
viscosity\footnote{It is important for convergence of the iterative solvers
 that the
``Laplace'' form of the Navier-Stokes equations is used.}
\begin{align}
 \rho \big( \frac{\partial \vek u}{\partial t} +
     \vek u \cdot \nabla \vek u \big) + \nabla p 
         - \eta \nabla^2 \vek u &=\vek f \label{eq:ns1} \\
\nabla \cdot \vek u &= 0
\label{eq:ns2}
\end{align} 
can be written in the following generic
form\footnote{In this chapter we will be using the notation of the book
by Elman \emph{et al.} \cite{Elman05} and the recent paper by Kay \emph{et al.}
\cite{Kay08}.}
\begin{equation}
  \begin{bmatrix}
    F & B^T \\ 
    B & 0
  \end{bmatrix}
  \begin{bmatrix}
     u \\ p
  \end{bmatrix}=
  \begin{bmatrix}
     f \\ g
  \end{bmatrix}
 \label{eq:systemns}
\end{equation}
where the matrix $F$ can be written as
\begin{equation}
   F = a M + b A + c N + W
\end{equation}
where $M$ is the mass matrix, $A$ is the discrete vector-diffusion (Laplacian)
matrix, $N$ is the vector convection matrix and
$W$ is the Newton derivative
matrix ($W=0$ for Picard iteration). All constant factors are put into 
the coefficients $a$, $b$ and $c$, for example $a=\rho/\Delta t$, $b=\eta$,
$c=\rho$ for an implicit Euler scheme. The first three terms of $F$ have a
block diagonal structure, where the number of diagonal blocks is equal to the
space dimension and each block is the same. If we denote this block 
as $F_v$, we find
\begin{equation}
   F_v = a M_v + b A_v + c N_v(\vek w)
\end{equation}
where the individual matrices are given by
\begin{align}
   M_v&=[M_v]_{ij}=(\phi_i,\phi_j), \\
   A_v&=[A_v]_{ij}=(\nabla\phi_i,\nabla\phi_j),  \\
   N_v(\vek w)&=[N_v]_{ij}=(\phi_i,\vek w\cdot\nabla\phi_j)
\end{align}
Here $\phi_i$ are the velocity shape functions and
$\vek w$ is a velocity vector depending on the time-integration and/or
iteration scheme used.

If we want to solve Eq.~\eqref{eq:systemns} using an iterative solver, we need 
a good preconditioner. A good preconditioner is found by the following block
LU-decomposition \cite{Elman05}:
\begin{equation}
  \begin{bmatrix}
    F & B^T \\ 
    B & 0
  \end{bmatrix}=
  \begin{bmatrix}
    I & 0 \\ 
    BF^{-1} & I
  \end{bmatrix}
  \begin{bmatrix}
    F & B^T \\ 
    0 & -BF^{-1}B^T
  \end{bmatrix}
\end{equation}
So if the (inverse of the) upper-triangular factor
\begin{equation}
   U=  \begin{bmatrix}
    F & B^T \\ 
    0 & -BF^{-1}B^T
  \end{bmatrix}
\end{equation}
is used as a preconditioner, i.e.\ the preconditioned matrix becomes
$L=KU^{-1}$ with $K$ the coefficient matrix in Eq.~\eqref{eq:systemns}, then
all the eigenvalues of the preconditioned matrix are identically one. The right
lower block in $U$ is called the Schur complement $S$:
\begin{equation}
   S=BF^{-1}B^T
\end{equation}
so that
\begin{equation}
   U=  \begin{bmatrix}
    F & B^T \\ 
    0 & -S
  \end{bmatrix}
\end{equation}
The preconditioning operation needed is $U^{-1}$ on a vector, which reduces to
solving the following system
\begin{equation}
  \begin{bmatrix}
    F & B^T \\ 
    0 & -S
  \end{bmatrix}
  \begin{bmatrix}
     w \\ s
  \end{bmatrix}=
  \begin{bmatrix}
     v \\ q
  \end{bmatrix}
\end{equation}
or
\begin{equation}
  S s = -q,\quad F w = v - B^T s
\end{equation}
Using $U^{-1}$ as a preconditioner combined with a GMRES iterative solver will
lead to the solution in just two iterations independent of the mesh and the
Reynolds number \cite{Elman05}. The last equation can be solved using a fast
iterative solver using an algebraic multigrid (AMG) preconditioner 
such as the HSL routine \texttt{hsl\_mi20}, see \cite{Doyle07}.
A big problem, however, is the construction of the Schur
complement matrix $S$, which is expensive and therefore
not practical. We need an approximation to the Schur complement
$\hat S$ such that the operation $\hat S^{-1}$ can be computed in an
inexpensive way. Since this will lead to more iterations, we also use an
inexpensive approximation for $F^{-1}$: one or two AMG V-cycles (as in
\cite{Kay08}).

A good approximation to the Schur complement is given by the so-called
pressure convection-diffusion preconditioner \cite{Kay08,Elman05}:
\begin{equation}
   \hat S = A_p F_p^{-1}M_p
\end{equation}
with the matrices $M_p$, $A_p$ and $F_p$ defined for the pressure system
in a similar way as for the velocity system:
\begin{align}
   F_p&= a M_p + b A_p + c N_p(\vek w), \\
   M_p&=[M_p]_{ij}=(\psi_i,\psi_j), \\
   A_p&=[A_p]_{ij}=(\nabla\psi_i,\nabla\psi_j),  \\
   N_p(\vek w)&=[N_p]_{ij}=(\psi_i,\vek w\cdot\nabla\psi_j)
\end{align}
Here $\psi_i$ are the pressure shape functions.
The operation $\hat S^{-1}$ can now be performed:
\begin{equation}
   \hat S^{-1} = M_p^{-1} F_pA_p^{-1}
   \label{eq:expressS}
\end{equation}
The inverse operation $A_p^{-1}$ will be \emph{approximated}
using one or two AMG
V-cycles (as in \cite{Kay08}) and $M_p^{-1}$ by several (typically eight)
diagonally (Jacobi) scaled conjugate gradient iterations. This has been
implemented in the routine \texttt{hsl\_solve3\_mi20} of the add-on
hsl\_extra (see Sec.~\ref{sec:hsl_extra_addon}).

For a Stokes system $a=c=0$, $b=\eta$ and therefore
 $F_v=\eta A_v$, $F_p=\eta A_p$ and
\begin{equation}
   \hat S = \frac1\eta M_p
\label{eq:stokes_precon}
\end{equation}
so $\hat S$ reduces to the (viscosity scaled) pressure mass matrix.
The inverse operation
$\hat S^{-1}$ can be approximated by several (typically eight)
diagonally (Jacobi) scaled conjugate gradient iterations.
This has been implemented in the routine \texttt{hsl\_solve1\_mi20}
of the add-on hsl\_extra (see Sec.~\ref{sec:hsl_extra_addon}).

For semi-implicit time-integration schemes like CN/AB2 and GK2 (see
Sec.~\ref{sec:navierstokes}) $c=0$ and therefore
 $F_p=aM_p+bA_p$ and
\begin{equation}
   \hat S^{-1} = M_p^{-1} F_pA_p^{-1}= bM_p^{-1}+aA_p^{-1}=M_1^{-1}+M_2^{-1}
\label{eq:stokes_helmholtz_precon}
\end{equation}
where the matrices are given by $M_1=M_p/b$ and $M_2=A_p/a$.
The inverse operation
$M_1^{-1}$ can be approximated by several (typically eight)
diagonally (Jacobi) scaled conjugate gradient iterations and 
$M_2^{-1}$ by one or two AMG V-cycles.
This has been implemented in the routine \texttt{hsl\_solve2\_mi20}
of the add-on hsl\_extra (see Sec.~\ref{sec:hsl_extra_addon}).

Notes:
\begin{itemize}
   \item See \cite{Elman05} and \cite{Kay08} for the theory and a full list of
proper references.
   \item The Stokes system is solved using unsymmetric preconditioning and an
iterative method for unsymmetric systems (BiCGSTAB or GMRES). In the book Elman
\emph{et al.} \cite{Elman05} a symmetric preconditioner (using the block
diagonal only) and an iterative method for indefinite symmetric systems
(MINRES) is
used. In the HSL-package MINRES is not available and we have not been able to
test whether this gives any advantages in CPU time.
\item The solvers require user input of the partitioning of the matrix
      and the pressure matrices (such $M_p$ etc.).
      The latter matrices require a separate pressure
      problem to be setup. Therefore, they are not as straightforward to use
      as direct solvers or ILU preconditioned iterative solvers. However, if
      the solvers are applicable for the problem at hand $\mathcal{O}(N)$
      scaling can be achieved, i.e.\ the CPU time and memory demands scale
      linearly with the number of degrees of freedom.
\item Currently the solvers are unable to handle constraints.
\item The viscous term needs to be build using the Laplace expression
      $\eta\nabla^2\vek u$ instead of $\nabla\cdot(2\eta\ten D)$ to get fast
      convergence. This affects the natural boundary conditions, but it also
remains uncertain whether generalized Newtonian systems can be solved. Maybe a
Picard scheme using $\nabla\cdot(\eta\nabla\vek u)$ can work.
\item For full Dirichlet conditions the pressure level can be left undetermined
      \cite{Elman05}.
\item There are some issues with the boundary conditions for the pressure
      system for inflow boundaries at high Reynolds number if convection is
      included in the pressure system ($N_p(\vek w)$) \cite{Elman05}.
\item The DEVSS-Stokes system can be handled, but requires more iterations than
      the Stokes system.
\end{itemize}

\subsection{Examples for Stokes and Stokes including a mass term}

The examples are based on the cavity flow for 2D (\texttt{stokes1.f90}) and
3D (\texttt{stokes11.f90}). The following examples are available:
\begin{description}
   \item{\texttt{stokes26}:} Block preconditioning using the pressure mass form
\eqref{eq:stokes_precon} for the 2D cavity.
   \item{\texttt{stokes26}:} Block preconditioning using the pressure mass form
\eqref{eq:stokes_precon} for the 3D cavity.
   \item{\texttt{stokes28}, \texttt{stokes29} and \texttt{stokes31}:}
The following Helmholtz-Stokes system is solved:
\begin{alignat}{3}
a\vek u-&\nabla^2\vec u+\nabla p & &=\vek f,&&\qquad\text{in }\Omega
 \label{eq:stokes+mass1}\\
   &\nabla\cdot\vek u& &=0&&\qquad\text{in }\Omega
 \label{eq:stokes+mass2}
\end{alignat}
for the 2D cavity.
This form appears when applying semi-implicit 
time-integration schemes like CN/AB2 and GK2 (see
Sec.~\ref{sec:navierstokes}).
In example \texttt{stokes28} this is solved using
block preconditioning with the pressure mass
form \eqref{eq:stokes_precon} as used for the Stokes equation.
As the value of $a$ is increased the number of iterations needed rapidly
increases.
In the examples \texttt{stokes29} and \texttt{stokes31} this is solved using
block preconditioning with the pressure mass and diffusion matrix, using
the form \eqref{eq:expressS} and \eqref{eq:stokes_helmholtz_precon},
respectively, for the approximate inverse of the Schur complement.
Rapid convergence is maintained for any value of $a$.
   \item{\texttt{stokes30}:} Block preconditioning using the pressure mass and
diffusion matrix with \eqref{eq:expressS} for solving
Eqs.~\eqref{eq:stokes+mass1}--\eqref{eq:stokes+mass2} in 
 the 3D cavity.
\end{description}

\subsection{Examples for Navier-Stokes}
The following examples are available:
\begin{description}
\item{\texttt{navier\_stokes12}:} Similar to \texttt{navier\_stokes1}
 (steady flow in a 2D driven cavity see Sec.~\ref{sec:nssteady2D}),
but now using pressure convection-diffusion
preconditioning in the form \eqref{eq:expressS} for the approximate inverse of the Schur complement. A lot of tweaking can
be done here for improving the performance (iterative solver, parameters in AMG, accuracy parameters, Picard versus Newton,
resolution of the mesh, steps in Reynolds number, \ldots ). At least a good resolution of the mesh is very important.
\item{\texttt{navier\_stokes13}:} Similar to \texttt{navier\_stokes7}
(manufactured solution in a 2D cavity see Sec.~\ref{sec:testproblemNS} using
AB2/CN time integration),
but now using pressure mass and diffusion
preconditioning in the form \eqref{eq:stokes_helmholtz_precon}
for the approximate inverse of the Schur complement. It turns out that
for this problem the solution is \emph{unstable}, whereas the direct solver
(\texttt{navier\_stokes7}) seems to be stable. It is unclear how this can
happen. Maybe, the AB2/CN method is only marginally stable and
errors produced by the iterative solver triggers the instability.
\item{\texttt{navier\_stokes14}:} Similar to \texttt{navier\_stokes11}
(manufactured solution in a 2D cavity see Sec.~\ref{sec:testproblemNS} using
GK2 time integration),
but now using pressure mass and diffusion
preconditioning in the form \eqref{eq:stokes_helmholtz_precon}
for the approximate inverse of the Schur complement. It turns out that
for this problem the solution is \emph{stable}, like the direct solver
(\texttt{navier\_stokes11}). 
 A lot of tweaking can
be done here for improving the performance and accuracy (iterative solver,
parameters in AMG, accuracy parameters, \ldots ). If the mass dominates (large
Reynolds, small time steps) the number of iterations is low, but increasing the
accuracy of the iterative solver (and thus more iterations) maybe necessary to
obtain an accurate solution. In any case, this looks like a nice candidate for
a fast Navier-Stokes solver comparable to pressure-correction, but without the
problems of pressure boundary conditions.
\end{description}
